\section{Related Work}

\subsection{Privacy-Preserving RAG}
Standard implementations of RAG process and store corporate documents in plaintext within centralized vector databases, creating severe targets for exfiltration and inversion attacks. Recent privacy-preserving RAG frameworks attempt to mitigate this by implementing structural data isolation. For example, \textbf{VertiSplitRAG} forces the vertical partitioning of document attributes before embedding, aiming to prevent joint attribute exposure. Similarly, \textbf{DPVoteRAG} applies differential privacy principles by sending disjoint data subsets to isolated LLM instances and securely aggregating the responses via voting. While these architectures enhance privacy, they inherently destroy the semantic layout of the knowledge base, leading to catastrophic degradations in retrieval utility.

\subsection{Document Clustering and Graph Partitioning}
To scale RAG over massive corpora, modern systems employ semantic clustering to speed up retrieval. \textbf{K-Means} clustering is the standard baseline, attempting to group semantically similar embeddings into identical shards. Advanced frameworks, such as \textbf{GraphRAG}, utilize community detection algorithms like \textbf{Leiden} modularity over K-Nearest Neighbor (K-NN) graphs to generate hierarchical representations of the data. Furthermore, \textbf{SPLIT-RAG} utilizes a greedy heuristic to partition semantic domains. While these techniques successfully maximize semantic coherence and minimize context fragmentation, our research demonstrates that they inadvertently maximize the risk of semantic reconstruction attacks, confirming a No-Free-Lunch theorem for RAG partitioning.
